\input _template

\setlanguage{CS}

\Title{Základy počítání úhlů}

\MathcompsLink{zaklady-pocitani-uhlu-1}

\Author{Patrik Bak}

\sec Úvod

Geometrie je nejvizuálnější část matematiky -- umíme si ji kreslit na papír, v~softwaru jako \Link[https://www.geogebra.org/]{GeoGebra} a~můžeme se těšit, když na vlastní oči vidíme, že věci v~ní fungují. Proč tomu tak je však může být záhadou. Abychom tato tajemství pochopili, potřebujeme kousek po kousku budovat základnu znalostí a~vizuální intuici.

Jedna z~nejzákladnějších technik je rozumět světu úhlů a~umět jejich vlastnosti s~úspěchem používat v~různých typech úloh. Jen málo zajímavých příkladů úhly nepoužívá. Teorie kolem úhlů přitom není složitá, je však důležité ji budovat pomalu a~detailně. Začneme s~nejjednoduššími vlastnostmi, které ani nevyžadují práci s~kružnicemi. Už ty se však dají s~úspěchem využít v~mnoha úlohách.

V~tomto materiálu se budeme věnovat úhlům, ale záměrně vynecháme úhly na kružnicích, kterým se patří věnovat samostatnou pozornost. Na chvíli zapomeňme, že kružnice existují.

Na závěr úvodu jedna menší poznámka o~konvenci. V~tomto materiálu uvidíme mnoho obrázků trojúhelníků, ve kterých je bod~$A$ nakreslen nahoře. Ukazuje se, že tento způsob kreslení je mezinárodní standard, zejména na poli matematické olympiády a~podobných soutěží se to tak zvykne dělat. V~Česku a~na~Slovensku je však ve školách stále zvykem kreslit~$C$ nahoře. V~některých zemích lze vidět i~$B$ nahoře, např. na Ukrajině. Z~hlediska řešení úloh je vždy dobré nakreslit si nahoru ten bod, u~kterého obrázek vypadá co nejsymetričtěji, snadněji v~něm jdou vidět další symetrické věci. V~tomto materiálu budou úlohy zadávány tak, že ta symetrie půjde vidět více při~$A$ nahoře \SlightSmile.

\sec Základy světa úhlů

V~této sekci si odvodíme ty nejjednodušší vlastnosti související s~úhly. Pevně věřím, že klíčem ke~zvládnutí geometrie je porozumění věcem od základů a~do hloubky, proto se těmto jednoduchým vlastnostem budeme věnovat více.

Základní vlastnosti úhlů, které běžně používáme skoro bez přemýšlení při řešení úloh:

\Highlight{
    \begitems \style i
    \i Vrcholové úhly jsou shodné:
       \Image{angles-vertical.pdf}
    \i Souhlasné úhly jsou shodné:
       \Image{angles-corresponding.pdf}
    \i Střídavé úhly jsou shodné:
       \Image{angles-alternate.pdf}
    \i Vedlejší úhly mají součet $180^\circ$:
       \Image{angles-supplementary.pdf}
   \enditems
}

Tyto vlastnosti spolu velmi přirozeně souvisí, menší cvičení k~zamyšlení:

\Exercise{}{
    Uvědomte si, že:
    \begitems \style a
    \i z~vlastnosti o~vedlejších úhlech vyplývá vlastnost vrcholových úhlů (a~naopak)
    \i z~kterékoli dvojice vlastností o~vrcholových, souhlasných, střídavých úhlech vyplývá ta třetí vlastnost
    \enditems
}{
    \textbf{(a)} Vedlejší úhel $\alpha'$ k~danému úhlu $\alpha$ má velikost $180^\circ - \alpha$. Úhel vedlejší k~$\alpha'$ má zase velikost $180^\circ - (180^\circ - \alpha) = \alpha$. Ten je však zároveň vrcholový úhel k~$\alpha$, tedy jsou vrcholové úhly shodné.

    \Image{angles-supplementary-derive-vertical.pdf}

    Naopak, předpokládejme vlastnost o~vrcholových úhlech a~označme po sobě jdoucí úhly v~průsečíku dvou přímek $\alpha$, $\beta$, $\gamma$, $\delta$. Ze shodnosti vrcholových úhlů $\alpha = \gamma$ a~$\beta = \delta$. Součet všech čtyř je plný úhel, tedy $\alpha + \beta + \gamma + \delta = 360^\circ$. Po dosazení dostáváme $2\alpha + 2\beta = 360^\circ$, čili $\alpha + \beta = 180^\circ$ -- to je právě vlastnost vedlejších úhlů.

    \Image{angles-vertical-derive-supplementary.pdf}

    \textbf{(b)} Všechny tři vlastnosti říkají, že jistá dvojice úhlů má stejnou velikost.

    Podívejme se na trojici úhlů $\alpha$, $\beta$, $\gamma$ z~obrázku a~na to, co o~jejich vztazích tvrdí jednotlivé vlastnosti. Pro dvojici $\alpha$, $\beta$ při horním průsečíku tvrdí vlastnost vrcholových úhlů $\alpha = \beta$. Pro dvojici $\alpha$, $\gamma$ tvrdí vlastnost souhlasných úhlů $\alpha = \gamma$. Konečně pro dvojici $\beta$, $\gamma$ tvrdí vlastnost střídavých úhlů $\beta = \gamma$.

    \Image{angles-corresponding-alternate-supplementary-connection.pdf}

    Tři vlastnosti tedy tvrdí tři rovnosti mezi týmiž třemi úhly, jen přes různé dvojice. Jakmile platí libovolné dvě, třetí plyne tranzitivitou: např. pokud víme (V) a~(S), tak $\alpha = \beta$ a~$\alpha = \gamma$, takže $\beta = \gamma$, čili (St). Ostatní dvojice analogicky.
}

Následující tvrzení všichni známe, umíte ho však dokázat?

\Theorem{}{
    Součet úhlů v~trojúhelníku je $180^\circ$.
}{
    Veďme přes vrchol~$A$ přímku rovnoběžnou s~$BC$. Protože $AB$ je příčka rovnoběžek, úhel $\beta$ při~$A$ a~úhel $\beta$ při~$B$ jsou střídavé, tedy stejné. Analogicky přes příčku $AC$ jsou úhly $\gamma$ při~$A$ a~při~$C$ střídavé. Úhly $\beta$, $\alpha$, $\gamma$ leží vedle sebe podél přímky přes~$A$, tedy $\alpha + \beta + \gamma = 180^\circ.$

    \Image{angles-triangle.pdf}
}

Toto tvrzení se rozhodně dá zobecnit. Víme například, že součet úhlů ve~čtyřúhelníku je $360^\circ$. Jak je to v~5-úhelníku? A~jak v~67-úhelníku? Odpovědí je následující tvrzení:

\Theorem{}{
    Nechť $n \ge 3$ je přirozené číslo. Součet úhlů v~konvexním $n$-úhelníku je roven $(n-2) \cdot 180^\circ$.
}{
    \NamedProof{Důkaz 1 (matematická indukce).} Základ $n=3$: součet úhlů trojúhelníku je $180^\circ = (3-2)\cdot 180^\circ$. Indukční krok: nechť $A_1 A_2 \ldots A_{n+1}$ je konvexní $(n+1)$-úhelník. Úhlopříčkou $A_1 A_n$ jej rozdělíme na trojúhelník $A_1 A_n A_{n+1}$ a~konvexní $n$-úhelník $A_1 A_2 \ldots A_n$: úhly při vrcholech $A_2, \ldots, A_{n-1}$ patří celé $n$-úhelníku, úhel při $A_{n+1}$ patří celý trojúhelníku a~úhly při $A_1$, $A_n$ se rozdělí mezi oba útvary tak, aby jejich části daly původní vnitřní úhly $(n+1)$-úhelníku. Součet úhlů $(n+1)$-úhelníku je tedy
    $$
    \underbrace{180^\circ}_{\text{trojúhelník}} + \underbrace{(n-2) \cdot 180^\circ}_{n\text{-úhelník}} = (n-1) \cdot 180^\circ = \bigl((n+1) - 2\bigr) \cdot 180^\circ.
    $$

    \Image{angles-polygon.pdf}

    \NamedProof{Důkaz 2 (vějířová triangulace).} Veďme z~vrcholu $A_1$ úhlopříčky do každého nesousedního vrcholu $A_3, A_4, \ldots, A_{n-1}$. Vznikne $n-2$ trojúhelníků pokrývajících vnitřek $n$-úhelníku bez překrytí. Při každém vrcholu $A_k$ se úhly sousedních trojúhelníků při $A_k$ přesně poskládají na vnitřní úhel $n$-úhelníku při $A_k$, takže součet úhlů všech trojúhelníků se rovná součtu vnitřních úhlů $n$-úhelníku, čili $(n-2)\cdot 180^\circ$.

    \Image{angles-polygon-fan.pdf}

    \Remark Důkaz~2 je \uv{rozbalený} důkaz~1: indukce odřezává trojúhelníky jeden po druhém, zde je vidíme všechny najednou.

    \NamedProof{Důkaz 3 (vnitřní bod).} Zvolme bod $O$ uvnitř $n$-úhelníku a~spojme ho s~každým vrcholem. Vznikne $n$ trojúhelníků s~celkovým součtem úhlů $n\cdot 180^\circ$. Úhly trojúhelníků při každém vrcholu $A_k$ spolu tvoří vnitřní úhel $n$-úhelníku při~$A_k$; úhly trojúhelníků při~$O$ spolu tvoří plný úhel $360^\circ$. Tedy součet vnitřních úhlů $n$-úhelníku je
    $$
    n\cdot 180^\circ - 360^\circ = (n-2)\cdot 180^\circ.
    $$

    \Image{angles-polygon-interior.pdf}

    \Remark{nekonvexní mnohoúhelníky} Vzorec $(n-2)\cdot 180^\circ$ platí i~pro nekonvexní mnohoúhelníky, pokud všechny vnitřní úhly měříme směrem dovnitř -- tedy i~úhly větší než $180^\circ$ počítáme jako takové.

    Vějířová triangulace z~jednoho vrcholu nemusí fungovat -- některé úhlopříčky mohou procházet mimo mnohoúhelník. Platí však, že každý jednoduchý $n$-úhelník se dá rozdělit na $n-2$ trojúhelníků (ne nutně vějířovitě), což důkaz zachrání. Důvodem je indukce přes trojúhelník tvořený nějakým vrcholem $V$ a~jeho sousedy $U$, $W$, který leží celý uvnitř našeho $n$-úhelníku. Odříznutím $V$ úhlopříčkou $UW$ dostaneme $(n-1)$-úhelník; opakováním tedy $n-2$ trojúhelníků. Dá se dokázat, že takový trojúhelník vždy existuje (dokonce dva), viz \Link[https://en.wikipedia.org/wiki/Two_ears_theorem]{věta o~dvou uších}.

    K~důkazu~3: funguje jen pokud existuje vnitřní bod $O$ viditelný ze všech vrcholů (tj. úsečky $OA_i$ leží celé uvnitř). Pro některé nekonvexní mnohoúhelníky takový bod~$O$ nemusí existovat (zkuste takový nakreslit) -- v~takovém případě se důkaz~3 nedá přímo zachránit a~je třeba sáhnout po triangulaci popsané v~předchozím odstavci.
}

Teď uveďme ještě dvě jednoduchá, ale užitečná pomocná tvrzení:

\Exercise{}{
    Dokažte, že součet úhlů na obrázku je roven $180^\circ$.

    \Image{angles-complementary-parallel-statement.pdf}
}{
    Označme $\alpha$ úhel svíraný příčkou $AB$ s~rovnoběžkami na straně $B$ (jako na obrázku). Souhlasné úhly $\alpha$ při~$A$ a~$\alpha$ při~$B$ jsou shodné. Při vrcholu $A$ jsou $\alpha$ a~$\beta$ vedlejší, tedy
    $$
    \beta + \alpha = 180^\circ.
    $$

    \Image{angles-complementary-parallel-solution.pdf}
}

\Exercise{}{
    Dokažte, že úhel s~otazníkem na obrázku je roven $\alpha+\beta$.

    \Image{angles-complementary-in-triangle.pdf}
}{
    Ze součtu úhlů v~trojúhelníku $|\angle ACB| = 180^\circ - \alpha - \beta$. Vnější úhel při $C$ je vedlejší k~$|\angle ACB|$:
    $$
    180^\circ - |\angle ACB| = \alpha + \beta.
    $$

    \Image{angles-complementary-in-triangle.pdf}
}

Poslední dvě cvičení možná vypadala velmi triviálně. Realita však je, že v~praktických úlohách jsou velmi užitečná -- když máme obrovský netriviální obrázek a~máme provést za sebou velký počet úhlových operací, tak je opravdu velmi výhodné si i~jen kousek té práce zjednodušit -- v~případě rovnoběžek nemusíme uvažovat pomocný souhlasný úhel a~v~případě trojúhelníku nemusíme pracovat s~$180^\circ$.

\sec Základy světa délek

Zatím jsme fungovali pouze ve světě úhlů a~vůbec jsme neřešili délky. Geometrie však začne být zajímavá, když tyto světy začneme propojovat. Základem pro nás bude {\em shodnost trojúhelníků}. Připomeňme si kritéria.

\Highlight{
    Věty o~shodnosti trojúhelníků:
    \begitems \style i
    \i Věta $sss$: dva trojúhelníky jsou shodné, pokud se shodují ve všech třech stranách.
    \i Věta $sus$: dva trojúhelníky jsou shodné, pokud se shodují ve dvou stranách a~úhlu, který tyto strany svírají.
    \i Věta $usu$: dva trojúhelníky jsou shodné, pokud se shodují v~jedné straně a~nějakých dvou úhlech.
    \i Věta $Ssu$: dva trojúhelníky jsou shodné, pokud se shodují ve dvou stranách a~úhlu naproti delší z~těchto stran.
    \enditems
}

Zdůrazněme, že u~věty $sus$ je opravdu důležité, že shodný úhel je ten svíraný oběma stranami -- bez tohoto předpokladu shodnost nemusí fungovat, viz obrázek. Na něm máme dva trojúhelníky $ABC$ a~$A'B'C'$, pro které platí $|BC|=|B'C'|$, $|AC|=|A'C'|$, a~$\beta=\beta'$, avšak zjevně nejsou shodné.

\Image{angles-sas-warning.pdf}

Věta $Ssu$ by byla naše záchrana -- zde ji však aplikovat neumíme, protože strana, naproti které leží úhel $\beta$, tedy $|AC|$, není nejdelší, neboť je zjevně kratší než $|BC|$ \SmileWithTear.

Než půjdeme dál, uvědomme si, co tyto věty znamenají. Z~mého pohledu je dobrý pohled na shodnost takový: chceme sestrojit trojúhelník, když máme dány nějaké tři jeho elementy -- budou všechny sestrojitelné trojúhelníky shodné?

\Example{}{
    Přesvědčte se o~platnosti věty $sss$ o~shodnosti.
}{
    Představme si, že máme dány tři úsečky délek $a$, $b$, $c$, ze kterých umíme sestrojit trojúhelník. Bez újmy na obecnosti předpokládejme, že $a=\max\{a,b,c\}$. Začněme konstruovat trojúhelník úsečkou~$BC$ délky~$a$. Následně  sestrojíme dvě kružnice: (a) kružnici se středem v~$B$ a~poloměrem $c$; (b) kružnici se středem v~$C$ a~poloměrem $b$. Díky $a \ge b$ resp. $a \ge c$ obě tyto kružnice protínají úsečku~$BC$. V~případě $a=b+c$ se protínají právě v~jednom bodě na ní (což neodpovídá trojúhelníku), v~případě $a>b+c$ se neprotínají vůbec, a~v~případě $a<b+c$ se protínají ve dvou bodech $A$ a~$A'$.

    \Image{angles-sss-proof.pdf}

    Na konci si uvědomme, že trojúhelníky $ABC$ a~$A'BC$ jsou však zřejmě vzájemně zrcadlovými obrazy podle $BC$. Také pokud bychom začali jinou úsečkou, tak vytvoříme pouze posunutí/otočení této konfigurace.

    \Remark Uvědomme si, že jsme vlastně po cestě dokázali~trojúhelníkovou nerovnost. Ta se obvykle formuluje jako: součet kterýchkoli dvou stran trojúhelníku je větší než třetí. My jsme dokázali, že součet dvou nejkratších je větší než nejdelší. To zřejmě stačí, protože součet nejdelší a~kterékoli další je větší než ta poslední.
}

Podobně umíme zdůvodnit další kritéria:

\Exercise{}{
    Přesvědčte se o~platnosti věty $sus$ o~shodnosti.
}{
    Nechť jsou dány strany $b = |AC|$, $c = |AB|$ a~úhel $|\angle BAC| = \alpha$ mezi nimi. Sestrojme vrchol~$A$ a~úsečku $AB$ délky~$c$ -- ta je dána až na shodné zobrazení. Bod $C$ musí ležet na polopřímce z~$A$ svírající s~$AB$ úhel~$\alpha$ a~zároveň ve vzdálenosti~$b$ od~$A$. Polopřímky svírající s~$AB$ úhel $\alpha$ jsou dvě (po jedné na každé straně $AB$): na každé leží právě jeden bod ve vzdálenosti $b$ od~$A$, tedy dostaneme body $C$ a~$C'$. Trojúhelníky $ABC$ a~$ABC'$ jsou zrcadlovými obrazy podle $AB$, tedy shodné.

    \Image{angles-sas-proof.pdf}
}

\Exercise{}{
    Přesvědčte se o~platnosti věty $usu$ o~shodnosti.
}{
    V~prvé řadě si uvědomme, že nezáleží na tom, které dva úhly jsou shodné -- pokud se dva trojúhelníky shodují ve dvou úhlech, tak třetí je určen jednoznačně, neboť všechny tři mají součet úhlů $180^\circ$.

    Pro účely našeho důkazu uvažme tedy, že jsou shodné právě úhly přiléhající ke shodné straně. Nechť je to úhel $\beta$ při vrcholu $B$, $\gamma$ při vrcholu $C$ a~strana $a = |BC|$ mezi nimi. Sestrojme úsečku $BC$ délky~$a$ -- ta je dána až na shodné zobrazení. Bod $A$ musí ležet na polopřímce z~$B$ svírající s~$BC$ úhel $\beta$ a~zároveň na polopřímce z~$C$ svírající s~$CB$ úhel $\gamma$. Součet vnitřních úhlů při $B$ a~$C$ v~trojúhelníku je menší než $180^\circ$, čili $\beta + \gamma < 180^\circ$, a~proto nejsou obě polopřímky rovnoběžné -- protnou se v~jediném bodě~$A$. Trojúhelník $ABC$ je tedy určen jednoznačně (až na shodná zobrazení).

    \Image{angles-asa-proof.pdf}
}

Méně známé tvrzení $Ssu$ je už trošičku těžší, proto k~němu budou i~návody \SlightSmile.

\Problem{0}{shodnost $Ssu$}{
    Přesvědčte se o~platnosti věty $Ssu$ o~shodnosti. Navíc si uvědomte, kde by se konstrukce pokazila, kdyby úhel nebyl naproti větší ze stran -- co by se stalo, kdyby byly obě strany stejně dlouhé?
}{
    Začněte konstrukci konstrukcí kratší úsečky. Jak pokračuje konstrukce po jejím sestrojení?
}{
    Po sestrojení kratší úsečky sestrojíme polopřímku pod naším daným shodným úhlem. Zbývá poslední krok konstrukce. Rozmyslete si, kdy dostaneme žádný, kdy jeden a~kdy dva vyhovující trojúhelníky.
}{
    Nechť jsou dány strany $a = |BC|$, $b = |AC|$ a~úhel $|\angle ABC| = \beta$ naproti delší straně $b$. Začneme sestrojením kratší úsečky $BC$ délky~$a$. Bod $A$ musí ležet na polopřímce z~$B$ svírající s~$BC$ úhel $\beta$ a~zároveň ve vzdálenosti $b$ od~$C$, čili na kružnici $k$ se středem~$C$ a~poloměrem $b$.

    Vzdálenost bodu~$B$ od středu kružnice je $|BC| = a$. Protože $a < b$, bod~$B$ leží \textit{uvnitř} kružnice $k$. Polopřímka z~bodu~$B$ tak začíná uvnitř kružnice a~pokračuje směrem ven, takže kružnici protne právě jednou. Bod $A$ je tedy určen jednoznačně, a~trojúhelník $ABC$ je rovněž jednoznačný (až na shodné zobrazení).

    \Image{angles-Ssa-proof-one-solution.pdf}

    Podívejme se, proč je podmínka $b > a$ podstatná. Pokud by bylo $a = b$, bod~$B$ by ležel přímo \textit{na} kružnici~$k$. Polopřímka z~$B$ by ji protínala v~bodě~$B$ samotném (což nedává trojúhelník) a~v~právě jednom dalším bodě -- ten je hledaným~$A$. Konstrukce tedy stále vede k~jednoznačnému (rovnoramennému) trojúhelníku, takže věta $Ssu$ formálně platí i~v~tomto hraničním případě. Pod $Ssu$ ji však samostatně nezařazujeme: pokud víme, že trojúhelník je rovnoramenný se $|BC|=|AC|$ a~známe jeden úhel $\beta$, tak ostatní dva úhly jsou už díky rovnoramennosti určeny (oba při základně mají velikost~$\beta$). Tatáž konfigurace je proto pokryta větami $sus$ či $sss$ a~$Ssu$ k~ní nepřidává žádnou novou informaci.

    \Image{angles-Ssa-proof-isosceles.pdf}

    Pokud by bylo $b < a$, bod $B$ by ležel \textit{vně} kružnice $k$. Polopřímka by ji mohla protnout ve dvou bodech -- vznikly by tak dva různé (neshodné) trojúhelníky.

    \Image{angles-Ssa-proof-two-solutions.pdf}

    \Remark Důkaz je možné provést i~tak, že nejprve sestrojíme delší úsečku. To však vyžaduje znalost množiny bodů nad pevnou úsečkou majících pevný úhel -- k~tomu se dostaneme v~jiném materiálu.
}

Na naší geometrické cestě se průběžně setkáme se všemi těmito tvrzeními. Prozatím si však ukažme nějaké konkrétní aplikace. Začneme tím nejzřejmějším tvrzením, které se ale také sluší a~patří dokázat:

\Theorem{rovnoramenný trojúhelník}{
    Dokažte, že pokud pro trojúhelník~$ABC$ platí $|AB|=|AC|$, tak $|\angle ABC|=|\angle ACB|$. Dokažte také obrácenou implikaci (tedy že z~rovnosti úhlů vyplývá rovnost délek).
}{
    \textit{Přímá implikace.} Předpokládejme $|AB| = |AC|$. Nechť $M$ je střed $BC$. Trojúhelníky $ABM$ a~$ACM$ jsou shodné podle $sss$ ($|AB|=|AC|$ z~předpokladu, $|BM|=|CM|$ ze~středu, společná strana $AM$), odkud $|\angle ABC| = |\angle ACB|$. Tento důkaz funguje i s~jinými volbami bodu $M$. Pokud za $M$ vezmeme patu kolmice z~$A$ na $BC$, dostaneme společnou stranu $AM$, $|AB|=|AC|$ a~pravý úhel při~$M$, který leží naproti nejdelší straně (přeponě) $AB$ resp.~$AC$ -- aplikujeme větu $Ssu$. Pokud za $M$ vezmeme patu osy vnitřního úhlu při~$A$, máme $|AB|=|AC|$, shodný úhel $|\angle BAM|=|\angle CAM|$ a~společnou stranu $AM$ mezi nimi -- aplikujeme větu $sus$.

    \Image{angles-isosceles-triangle.pdf}

    \textit{Obrácená implikace.} Předpokládejme $|\angle ABC| = |\angle ACB| = \beta$. Nechť $M$ je pata kolmice z~$A$ na $BC$. Podle věty $usu$ je $\triangle ABM \cong \triangle ACM$, neboť máme dva shodné úhly $|\angle ABM| = |\angle ACM|$ a~$|\angle AMB| = |\angle AMC| = 90^\circ$ a~společnou stranu~$AM$, tedy $|AB| = |AC|$. Podobně bychom za $M$ mohli zvolit patu osy vnitřního úhlu při~$A$ -- ze shodných úhlů $|\angle BAM|=|\angle CAM|$ a~$|\angle ABM|=|\angle ACM|$ a~ze společné strany $AM$ znovu aplikujeme větu $usu$. Volba $M$ jako středu $BC$ tu však nepomáhá: dostaneme sice $|BM|=|CM|$, $AM$ společnou a~$|\angle ABM|=|\angle ACM|=\beta$, ale úhel $\beta$ leží naproti $AM$, která může být kratší než $BM$ (při tupém úhlu při~$A$), takže ani větu $Ssu$ obecně použít neumíme -- ponaučení je, že i~když máme správný bod, na jeho přesné definici záleží (to ještě mnohokrát uvidíme).

    \Remark Jiný vtipný důkaz je založen na tom, že dokážeme $\triangle ABC \cong \triangle ACB$ (všimněte si různého pořadí vrcholů). V~případě, že známe stejné strany, použijeme větu $sus$ nebo dokonce $sss$ -- ta nám následně dá shodné úhly. V~případě, že známe stejné úhly, to zase bude věta $usu$, a~shodnost nám dá stejné strany.
}

Vyzkoušejte si exaktně dokázat tyto jednoduché poznatky, to už bude jednodušší:

\Exercise{}{
    Rovnostranný trojúhelník má tři shodné úhly rovny $60^\circ$.
}{
    V~$\triangle ABC$ s~$|AB| = |AC| = |BC|$ dává věta o~rovnoramenném trojúhelníku $|\angle ABC| = |\angle ACB|$ (z~$|AB|=|AC|$); analogicky z~$|AB|=|BC|$ je $|\angle BAC| = |\angle BCA|$. Všechny tři úhly jsou tedy shodné, a~ze součtu $180^\circ$ je každý $60^\circ$.

    \Image{angles-equilateral.pdf}
}

\Exercise{}{
    Rovnoramenný pravoúhlý trojúhelník má tři úhly rovny $90^\circ$, $45^\circ$, $45^\circ$.
}{
    Nechť je pravý úhel při vrcholu~$A$, tedy $|\angle BAC| = 90^\circ$, a~odvěsny jsou shodné, $|AB| = |AC|$. Podle věty o~rovnoramenném trojúhelníku platí $|\angle ABC| = |\angle ACB|$. Označme tuto společnou velikost $\beta$. Ze součtu úhlů v~trojúhelníku
    $$
    90^\circ + \beta + \beta = 180^\circ,
    $$
    odkud $2\beta = 90^\circ$, čili $\beta = 45^\circ$.

    \Image{angles-isosceles-right-triangle.pdf}
}

Velmi šikovná věc použitelná v~nelehkých úlohách je následující úloha. Jeden možný důkaz je přes trigonometrii. My to ale chceme pěkně geometricky.

\Problem{0}{}{
    Pravoúhlý trojúhelník má zbývající dva úhly rovny $60^\circ$ a~$30^\circ$ právě tehdy, když je jeho přepona dvakrát delší než kratší odvěsna.
}{
    Řekněme, že $AB$ je přepona a~je dvakrát delší než odvěsna $BC$. Trikem je uvážit bod $B'$ takový, že $C$ je střed úsečky $BB'$.
}{
    Nechť náš trojúhelník má pravý úhel při vrcholu~$C$.

    \textit{($\Rightarrow$)} Předpokládejme $|AB| = 2|BC|$. Nechť $B'$ je takový bod, že $C$ je středem úsečky $BB'$. Potom $|B'C| = |BC|$ a~$C$ leží mezi $B$, $B'$, takže $|\angle B'CA|$ je vedlejší k~$|\angle BCA| = 90^\circ$, čili také $90^\circ$. Platí $\triangle CBA \cong \triangle CB'A$ z~věty $sus$, neboť úhly při~$C$ jsou oba pravé, strana~$AC$ je společná, a~$|BC|=|B'C|$.

    \Image{angles-90-60-30.pdf}

    Protože $C$ je středem $BB'$, je $|BB'| = 2|BC|$, a~ze zadání $|AB| = 2|BC|$. Spolu
    $$
    |BB'| = |AB| = |AB'|,
    $$
    takže $\triangle ABB'$ je rovnostranný a~všechny jeho úhly jsou $60^\circ$. Speciálně $|\angle ABC| = |\angle ABB'| = 60^\circ$. Ze součtu úhlů v~$\triangle ABC$ dopočítáme $|\angle BAC| = 180^\circ - 90^\circ - 60^\circ = 30^\circ$.

    \textit{($\Leftarrow$)} Není těžké rozmyslet si, že úvahy z~předchozího odstavce umíme snadno obrátit -- klíčovou shodnost $\triangle CBA \cong \triangle CB'A$ tentokrát dostaneme z~věty $usu$.

    \Remark Jiné řešení je uvážit, že podle Thaletovy věty je střed~$O$ kružnice opsané trojúhelníku $ABC$ zároveň středem přepony $AB$. K~tomuto řešení se podrobněji vrátíme, až se budeme bavit o~kružnicích.
}

\sec Co si zapamatovat

\secc Techniky

\begitems \style N
\i Při rovnoběžkách hledáme shodné střídavé a~souhlasné úhly, případně úhly se součtem~$180^\circ$.
\i Je užitečné nahlížet na vnější úhel trojúhelníku jako na součet dvou úhlů.
\i Pro propojení světa délek a~úhlů se hodí shodnost dvou trojúhelníků.
\i Pomocné body (střed úsečky, zrcadlový obraz, prodloužení) často odhalí shodnost nebo speciální trojúhelník.
\enditems

\secc Užitečná fakta

\begitems \style N
\i Součet úhlů v~trojúhelníku je $180^\circ$, obecně v~konvexním $n$-úhelníku $(n-2) \cdot 180^\circ$.
\i Věty o~shodnosti trojúhelníků: $sss$, $sus$, $usu$, $Ssu$.
\i V~trojúhelníku jsou dvě strany shodné právě když jsou shodné jim protilehlé úhly.
\i Pravoúhlý trojúhelník má přeponu dvakrát delší než odvěsnu právě když má úhly $90^\circ, 60^\circ, 30^\circ$.
\enditems

\sec Úlohy

V~této sekci najdete různé úlohy, na které vám stačí poznatky z~tohoto materiálu, žádné pokročilé věci jako \textit{obvodové a~středové úhly} potřeba nejsou (na tyto úlohy se podíváme později).


\Problem{0}{MO okresní kolo Z8 2025}{
    V~trojúhelníku $ABC$ leží bod $D$ na straně $BC$, bod $E$ na straně $AC$, přičemž $|AB| = |BE| = |EC| = |CD|$ a~$|BD| = |DE|$. Určete velikosti úhlů $ACB$ a~$BAD$.
}{
    Označte $\gamma = |\angle ACB|$ a~snažte se všechny úhly v~obrázku vyjádřit pomocí~$\gamma$. Máme spoustu rovnoramenností.
}{
    Jedna možná cesta je postupně pomocí $\gamma$ vyjádřit úhly $CBE$, $BED$, $CDE$, $CED$. Neumíme teď někde napsat rovnici, ze které se dá vypočítat $\gamma$?
}{
    Podívejte se na součet úhlů v~trojúhelníku~$BEC$. Z~něj by mělo jít spočítat $\gamma$.
}{
    K~dopočítání~$BAD$ potřebujeme více kroků. První je spočítat co nejvíce úhlů a~něčeho si všimnout.
}{
    Když spočítáme úhly v~$ABC$, dostaneme, že $ABC$ je rovnoramenný.
}{
    Teď už umíme použít symetrii, jedna možnost je dokázat shodnost trojúhelníků $EAD$ a~$DBE$.
}{
    Označme $\gamma = |\angle ACB|$.

    V~rovnoramenném $\triangle BEC$ ($|BE|=|EC|$) je $|\angle EBC| = |\angle ECB| = \gamma$. Protože $D$ leží na $BC$, je $|\angle DBE| = |\angle EBC| = \gamma$, a~z~rovnoramennosti $\triangle BDE$ ($|BD|=|DE|$) také $|\angle DEB| = \gamma$. Věta o~vnějším úhlu v~$\triangle BDE$ při~$D$ dává
    $$
    |\angle EDC| = |\angle DBE| + |\angle DEB| = 2\gamma,
    $$
    a~v~rovnoramenném $\triangle CDE$ ($|CD|=|CE|$) tak $|\angle CED| = |\angle CDE| = 2\gamma$. Ze součtu úhlů v~něm $\gamma + 4\gamma = 180^\circ$, čili $\gamma = 36^\circ$.

    \Image{angles-isosceles-chain.pdf}{0.8}

    Nyní si vezměme, že $|\angle AEB|$ je vedlejší úhel k~$|\angle BEC|=3\gamma=108^\circ$, takže je roven $72^\circ$. Z~rovnoramennosti $\triangle BAE$ tak také $|\angle BAE|=72^\circ$. Z~$\triangle ABC$, ve kterém už máme při $C$ a~$A$ po řadě $36^\circ$ a~$72^\circ$, dopočítáme $|\angle CBA|=72^\circ$, takže trojúhelník $ABC$ je rovnoramenný.

    Dále z~$|CA|=|CB|$ a $|CE|=|CD|$ vlastně máme $|EA|=|DB|$. Vezměme trojúhelníky $EAD$ a $DBE$; ty jsou rovnoramenné se shodným úhlem při vrcholu (rovným $180^\circ-2\gamma$) a~shodnými rameny, takže podle $sus$ jsou shodné. To dává $|\angle DAE|=|\angle DBE|=36^\circ$.

    Na závěr tedy máme
    $$
    |\angle BAD| = |\angle BAE| - |\angle DAE| = 2\gamma - \gamma = \gamma = 36^\circ.
    $$
}

\Problem{0}{DuoGeo 2025}{
    Do čtverce $ABCD$ byly nakresleny rovnostranné trojúhelníky $ABX$ a~$CDY$. Určete součet vyznačených úhlů.

    \Image{angles-square-equilateral-statement.pdf}{0.7}
}{
    Zkuste nejprve spočítat co nejvíce úhlů na obrázku.
}{
    Klíčem k~dopočítání finálního úhlu je najít vhodný rovnoramenný trojúhelník přes stejné délky.
}{
    Podle symetrie jsou všechny čtyři vyznačené úhly shodné. Stačí tedy najít velikost jednoho z~nich -- zaměříme se na $|\angle YDX|$.

    Nejprve $|\angle XAD| = |\angle BAD| - |\angle BAX| = 90^\circ - 60^\circ = 30^\circ$; analogicky $|\angle ADY| = |\angle ADC| - |\angle YDC| = 90^\circ - 60^\circ = 30^\circ$.

    Dále $|DA| = |AB| = |AX|$ (první je strana čtverce, druhé strana rovnostranného trojúhelníka $ABX$), takže $\triangle AXD$ je rovnoramenný se základnou $DX$. Z~rovnosti $|\angle XAD| = 30^\circ$ a~ze součtu úhlů v~trojúhelníku
    $$
    |\angle ADX| = \tfrac{1}{2}\bigl(180^\circ - 30^\circ\bigr) = 75^\circ.
    $$
    Hledaný úhel je rozdílem dvou již vypočítaných:
    $$
    |\angle YDX| = |\angle ADX| - |\angle ADY| = 75^\circ - 30^\circ = 45^\circ.
    $$
    Součet všech čtyř vyznačených úhlů je tedy $4 \cdot 45^\circ = 180^\circ$.

    \Image{angles-square-equilateral-solution.pdf}{0.8}
}

\Problem{0}{MO školní kolo C 2024}{
    V~lichoběžníku $ABCD$, kde $AB \parallel CD$, se osy vnitřních úhlů při vrcholech $C$ a~$D$ protínají na úsečce~$AB$. Dokažte, že $|AD| + |BC| = |AB|$.
}{
    Nechť~$P$ je náš společný průsečík. Máme rovnoběžnost, máme osu úhlu, to dává dost stejných úhlů.
}{
    Cílem je najít rovnoramenné trojúhelníky.
}{
    Označme $P$ společný průsečík os; podle zadání leží na úsečce~$AB$. Ukážeme, že
    $$
    |AD| = |AP| \quad \hbox{a} \quad |BC| = |BP|,
    $$
    odkud přímo $|AD| + |BC| = |AP| + |BP| = |AB|$.

    Protože $DP$ je osa vnitřního úhlu při~$D$, máme $|\angle ADP| = |\angle CDP|$. Z~rovnoběžnosti $AB \parallel CD$ jsou $|\angle CDP|$ a~$|\angle APD|$ střídavé úhly při příčce~$DP$, takže
    $$
    |\angle APD| = |\angle CDP| = |\angle ADP|.
    $$
    V~trojúhelníku $ADP$ jsou tedy úhly při vrcholech $D$ a~$P$ shodné, takže $|AD| = |AP|$. Analogickou úvahou při vrcholu~$C$ dostaneme $|BC| = |BP|$.

    \Image{angles-trapezoid-bisectors.pdf}{0.8}
}

\Problem{0}{MO školní kolo A 2023}{
    V~konvexním pětiúhelníku $ABCDE$ platí $|\angle CBA| = |\angle BAE| = |\angle AED|$. Na stranách $AB$ a~$AE$ existují po řadě body $P$ a~$Q$ tak, že $|AP| = |BC| = |QE|$ a~$|AQ| = |BP| = |DE|$. Dokažte, že $CD \parallel PQ$.
}{
    Máme spoustu stejných úhlů a~stran, zkuste najít shodné trojúhelníky.
}{
    Klíčové jsou shodné trojúhelníky $PBC$, $QAP$ a~$DEQ$. Z~nich dostaneme užitečné shodné vlastnosti. Pokračujte v~hledání shodností.
}{
    Finální shodnost k~dokázání je $CPQ$ a $DQP$. Proč to stačí?
}{
    Protože $|BC| = |AP| = |EQ|$, $|BP| = |AQ| = |ED|$ a~$|\angle CBP| = |\angle PAQ| = |\angle QED|$, jsou podle věty $sus$ trojúhelníky $PBC$, $QAP$ a~$DEQ$ navzájem shodné.

    \Image{angles-pentagon.pdf}{0.8}

    Odtud plyne $|CP| = |PQ| = |QD|$ a~také
    $$
    \gather
    |\angle CPQ| = 180^\circ - |\angle BPC| - |\angle APQ| = \\ = 180^\circ - |\angle PQA| - |\angle EQD| = |\angle PQD|.
    \endgather
    $$
    Podle věty $sus$ jsou tedy shodné i~rovnoramenné trojúhelníky $CPQ$ a~$DQP$. Z~toho vyplývá shodnost jejich výšek z~vrcholů $C$ a~$D$ na společnou stranu $PQ$; tyto výšky jsou navíc rovnoběžné (obě kolmé na $PQ$), takže $CD \parallel PQ$.
}

\Problem{0}{DuoGeo 2025}{
    Je dán čtyřúhelník $ABCD$ s~průsečíkem úhlopříček $T$. Předpokládejme, že velikosti úhlů $BAC$ a~$DBA$ jsou po řadě $30^\circ$ a~$45^\circ$. Na úsečce $BT$ leží bod $Z$ takový, že $CZ \perp BT$. Předpokládejme, že přímka $CZ$ protne úsečku $AB$ v~bodě $M$. Nechť $R$ je průsečík úseček $AT$ a~$MD$. Předpokládejme, že $|AM| = |AR|$ a~$|MR| + |TD| = 14$. Určete velikost úsečky $|BZ|$.
}{
    Dopočítávejte úhly, dokud nenajdeme rovnoramenný trojúhelník.
}{
    Trpělivým počítáním úhlů se dá dojít k~tomu, že $DTR$ je rovnoramenný. To by mělo vnést světlo do podmínky $|MR|+|TD|=14$.
}{
    Podmínka se po dokázání rovnoramenností přeloží jako $|MD|=14$. To pomůže později -- aktuálně už víc úhlů nespočítáme a~musíme ještě něco najít ze světa délek. Klíčem je najít pěkný pravoúhlý trojúhelník.
}{
    Dá se spočítat, že úhly trojúhelníku $MDZ$ jsou $90-60-30$. Známe jeho přeponu $MD$. Naše skvěle dokázané pomocné tvrzení nám teď dává další délku. Pak je to krok od řešení.
}{
    V~trojúhelníku $ATB$ známe úhly při vrcholech $A$ a~$B$, a~sice $30^\circ$ a~$45^\circ$. Tím pádem je vnější úhel při~$T$ roven součtu těchto úhlů, konkrétně
    $$
    |\angle DTR|=|\angle TAB|=|\angle TBA|=30^\circ+45^\circ=75^\circ.
    $$
    V~rovnoramenném trojúhelníku $AMR$ s~$|AM| = |AR|$ je úhel při vrcholu $A$ roven $30^\circ$, takže oba úhly při základně $MR$ mají velikost $\tfrac{1}{2}(180^\circ - 30^\circ) = 75^\circ$. Speciálně $|\angle ARM| = 75^\circ$, a~tedy i~jeho vrcholový úhel $|\angle DRT|$ má velikost $75^\circ$.

    \Image{angles-quad-30-45.pdf}{0.8}

    Trojúhelník $DTR$ má tedy dva úhly velikosti $75^\circ$, čili je rovnoramenný se~základnou $TR$ a~$|DT| = |DR|$. Odtud
    $$
    |MD| = |MR| + |RD| = |MR| + |TD| = 14.
    $$
    Navíc $|\angle TDR| = 180^\circ - 2 \cdot 75^\circ = 30^\circ$.

    Podívejme se nyní na trojúhelník $MDZ$. Protože $T$, $Z$ leží na úsečce $BD$ a~$R$ leží na úsečce $MD$, je $|\angle MDZ| = |\angle RDT| = 30^\circ$. Trojúhelník $MDZ$ je tedy pravoúhlý (při~$Z$) s~přeponou $MD$ a~úhlem $30^\circ$ při~$D$, takže podle dřívějšího tvrzení o~$30^\circ$-$60^\circ$-$90^\circ$ trojúhelníku
    $$
    |MZ| = \tfrac{1}{2} |MD| = 7.
    $$

    Konečně trojúhelník $MZB$ má pravý úhel při~$Z$ a~úhel $45^\circ$ při~$B$, takže i~$|\angle BMZ| = 45^\circ$. Je tedy rovnoramenný pravoúhlý a~$|BZ| = |MZ| = 7$.
}

\DisplayExerciseSolutions

\DisplayHints

\DisplayProblemSolutions

\bye
